\documentclass[11pt]{article}
\usepackage[T1]{fontenc}
\usepackage{amsmath,amssymb,amsthm}
\usepackage[margin=1in]{geometry}
\usepackage{hyperref}
\newtheorem{theorem}{Theorem}
\newtheorem{lemma}[theorem]{Lemma}
\newtheorem{corollary}[theorem]{Corollary}
\newtheorem{definition}[theorem]{Definition}
\newcommand{\N}{\mathbb N}
\newcommand{\Z}{\mathbb Z}
\title{From One Robust Four-Chromatic Seed to All Sufficiently Large Orders}
\author{Research draft for independent verification}
\date{6 November 2026}
\begin{document}
\maketitle

\begin{abstract}
We prove an order-completion theorem: if a vertex-critical four-chromatic
graph on $g$ vertices remains four-chromatic after deletion of at most $r$
edges, then such a graph exists on every number $N\geq45g^2$ of vertices.
The proof combines the published substitution lemma of Skottov\'a and
Steiner with Haj\'os joins of three small critical graphs and an elementary
two-generator conductor bound. Consequently, the seed-existence theorem
of Kruer and Kohlmeyer implies $f_4(N)\to\infty$, the all-orders form of
Erd\H{o}s problem 944. The seed theorem is the external input; it is not
re-proved here. This manuscript makes no claim of independently established
publication priority for the order-completion argument.
\end{abstract}

\section{Statement and source input}

All graphs are finite and simple. A graph is \emph{vertex-critical
four-chromatic} if $\chi(G)=4$ and $\chi(G-v)=3$ for every vertex $v$.
We call $G$ a \emph{$(4,r)$-graph} if it is vertex-critical four-chromatic
and $\chi(G-F)=4$ for every set $F\subseteq E(G)$ with $|F|\leq r$.
In contrast, a \emph{4-critical skeleton} below is four-chromatic and
becomes three-colorable upon deletion of any edge or any vertex.
The distinction is essential: a $(4,r)$-graph for $r\geq1$ is not
edge-critical.

Following Skottov\'a and Steiner~\cite{SS}, let $f_4(N)$ be the largest
integer $r\geq1$ for which a $(4,r)$-graph of order $N$ exists, and set it
to zero if no such $r$ exists. Thus $r$ denotes the edge-deletion tolerance,
not the minimum size of an edge set whose deletion lowers the chromatic
number; the latter is at least $r+1$.

Kruer and Kohlmeyer~\cite[Theorem 1.1]{KK} assert that for every $r\geq1$
there exists a finite $(4,r)$-graph. Their introductory scope statement
does not claim existence at every sufficiently large order. The following
transfer supplies that implication from an arbitrary single seed.

\begin{theorem}[Single-seed order completion]\label{thm:main}
Let $r\geq1$, and suppose that a $(4,r)$-graph of order $g$ exists.
Then a $(4,r)$-graph exists on every integer number $N\geq45g^2$ of
vertices. More precisely, it suffices that
\[
N\geq
\begin{cases}
45g^2-46g+13,&g\text{ odd},\\
\dfrac{45}{2}g^2-8g+1,&g\text{ even}.
\end{cases}
\]
\end{theorem}

\begin{corollary}\label{cor:limit}
The seed-existence theorem of Kruer and Kohlmeyer implies
$\displaystyle\lim_{N\to\infty}f_4(N)=\infty$.
\end{corollary}

\begin{proof}[Proof of Corollary~\ref{cor:limit}]
Fix an arbitrary integer $r\geq1$ and choose a seed of order $g=g(r)$
using~\cite[Theorem 1.1]{KK}. Theorem~\ref{thm:main} gives $f_4(N)\geq r$
for every $N\geq45g(r)^2$. This is exactly the asserted divergence.
\end{proof}

\section{The published substitution lemma, with proof}

For a map $\psi:V(G)\to\{1,2,3\}$, write $m_G(\psi)$ for the number
of edges whose endpoints receive the same color. The edge-robustness
condition is equivalent to
\begin{equation}\label{eq:mono}
m_G(\psi)\geq r+1\qquad\hbox{for every }\psi:V(G)\to\{1,2,3\}.
\end{equation}
Indeed, deleting all monochromatic edges produces a proper three-coloring;
conversely, a proper three-coloring after deletion of at most $r$ edges
has at most $r$ monochromatic edges in the original graph.

\begin{lemma}[Skottov\'a--Steiner substitution]\label{lem:sub}
Suppose that a $(4,r)$-graph $G$ has $g$ vertices and that $H$ is a
4-critical skeleton with $h$ vertices and $e$ edges. Then there is a
$(4,r)$-graph of order $h+(g-1)e$.
\end{lemma}

This is the equal-seed case of~\cite[Lemma 2.1]{SS}. The proof is included
to make every criticality condition in the transfer explicit; the
substitution itself is not claimed as a new construction.

\begin{proof}
Choose $v\in V(G)$ and a proper three-coloring $\varphi$ of $G-v$.
Split $v$ into two nonadjacent terminals $a,b$. Attach $a$ to the
neighbors of $v$ of $\varphi$-color 1, and attach $b$ to the remaining
neighbors of $v$. Keep all edges of $G-v$, obtaining a graph $S$ with
$g-1$ internal vertices. We use the following properties.

First, $S$ has a proper three-coloring with distinct terminal colors:
keep $\varphi$ and assign $a=2,b=1$. Permuting the colors realizes any
prescribed pair of distinct terminal colors. Second, every three-color
assignment with equal terminal colors has at least $r+1$ monochromatic
edges: identifying the terminals recovers $G$ edge for edge, so
\eqref{eq:mono} applies. Third, deleting any internal vertex $x$ permits
a proper three-coloring with equal terminal colors, by taking a
three-coloring of $G-x$ and giving both terminals the color of $v$.
Permuting colors realizes any prescribed common color. Fourth, deleting
either terminal leaves a three-colorable graph, with any prescribed color
at the surviving terminal, by restricting and permuting the first coloring.

For every edge $xy$ of $H$, delete that edge and insert a fresh copy of
$S$, identifying its terminals with $x,y$. All gadget interiors are
disjoint. Call the resulting simple graph $T(H)$.

A proper four-coloring of $H$ extends to $T(H)$: each gadget has distinct
terminal colors, and its proper three-coloring can be relabeled using
those two colors and any third one from the four-color palette.
On the other hand, any three-color assignment to $T(H)$ restricts to an
assignment on $H$ with a monochromatic edge $xy$. The gadget replacing
that edge contributes at least $r+1$ monochromatic edges. Therefore no
deletion of at most $r$ edges makes $T(H)$ three-colorable, and
$\chi(T(H)-F)=4$ for all such $F$.

It remains to check every vertex deletion. If $x$ is internal to the
gadget replacing $uv$, take a proper three-coloring of $H-uv$.
The endpoints $u,v$ necessarily have equal colors, for otherwise this
would properly three-color $H$. Extend over the punctured gadget using
the third property, and over every other gadget using the first.
If instead a skeleton vertex $u$ is deleted, properly three-color $H-u$.
Use the first property for gadgets not incident with $u$ and the fourth
for each gadget that has lost its terminal $u$. Thus every vertex
deletion is three-colorable. Its chromatic number is exactly three,
since adding one vertex to a two-colorable graph cannot produce a
four-chromatic graph. Finally, each skeleton edge contributes exactly
$g-1$ new vertices, proving the order formula.
\end{proof}

\section{Three critical skeletons and Haj\'os joins}

We use skeletons with the following vertex and edge counts:
\[
\begin{array}{c|r|r}
\text{graph}&|V|&|E|\\\hline
K_4&4&6\\
C_5\text{ with an additional universal hub}&6&10\\
M(C_5)\text{, the Gr\"otzsch graph}&11&20
\end{array}
\]
All three are 4-critical skeletons. This is immediate for $K_4$.
The odd wheel is four-chromatic because its rim needs three colors and
its hub needs another. Removing a rim edge allows two colors on the
remaining rim path and a third at the hub. Removing a spoke allows the
hub and that spoke's rim endpoint to share a color; the remaining four
rim vertices form a path and can use the other two colors.
Every vertex deletion is then three-colorable: choose an incident edge
and restrict a three-coloring after deleting that edge.

For completeness, define $M(C_5)$ to have vertices
$v_0,\ldots,v_4,u_0,\ldots,u_4,w$, with subscripts modulo five, and
edges $v_iv_{i+1}$, $u_iv_{i-1}$, $u_iv_{i+1}$, and $wu_i$.
In a hypothetical three-coloring give $w$ color 3; all $u_i$ then have
colors 1 or 2. Recolor each $v_i$ of color 3 with the color of $u_i$.
Two adjacent $v_i$ cannot both have originally had color 3, and the
cross edges ensure that recoloring only one endpoint preserves properness.
This would give a proper two-coloring of $C_5$, a contradiction.
A four-coloring is obtained by three-coloring the rim, giving each $u_i$
the color of $v_i$, and giving $w$ a fourth color.

The dihedral symmetries have three edge orbits. Here are proper
three-colorings after deleting one representative of each, in vertex
order $(v_0,\ldots,v_4,u_0,\ldots,u_4,w)$, using colors $0,1,2$:
\[
\begin{array}{c|c}
\text{deleted edge}&\text{color tuple}\\\hline
v_0v_1&(1,1,0,1,0,\,2,2,2,1,2,\,0)\\
u_0v_1&(0,1,0,1,2,\,1,1,2,1,2,\,0)\\
wu_0&(0,1,0,1,2,\,0,1,2,1,2,\,0)
\end{array}
\]
Inspection of the displayed edge definition verifies each tuple.
Symmetry covers every edge. There are no isolated vertices, so
edge-deletion colorability also gives vertex-deletion colorability.

\begin{lemma}[Haj\'os join]\label{lem:hajos}
If $H_1,H_2$ are 4-critical skeletons, there is a 4-critical skeleton $J$
with
\[
|V(J)|=|V(H_1)|+|V(H_2)|-1,\qquad
|E(J)|=|E(H_1)|+|E(H_2)|-1.
\]
\end{lemma}

\begin{proof}
Start with disjoint copies. Choose edges $a_1b_1$ and $a_2b_2$, delete
them, identify $a_1,a_2$ to one vertex $a$, and insert edge $b_1b_2$.
The counts follow immediately, and the graph remains simple.
In any proper three-coloring of $H_i-a_ib_i$, the distinguished endpoints
must have equal colors. Thus a proper three-coloring of $J$ would force
$b_1,b_2$ to agree, contrary to their new edge. A four-coloring exists:
take proper four-colorings of the two summands, permute them to match
at $a$, and permute the other three colors on one summand so that
$b_1,b_2$ differ.

Deleting the new edge permits three-colorings of both edge-deleted
summands, matched at $a$. If an old edge $f$ from the first summand is
deleted, take a three-coloring of $H_1-f$. Its distinguished endpoints
$a_1,b_1$ differ, because their original edge is present in that graph.
Combine it with a three-coloring of $H_2-a_2b_2$, whose distinguished
endpoints agree, matching the colors at $a$. The new edge is proper.
The other summand is symmetric. Hence every edge deletion is
three-colorable. Every vertex has an incident edge, so restricting
one of these colorings proves vertex-deletion colorability as well.
\end{proof}

\section{Order completion and the explicit threshold}

\begin{lemma}[Two-generator conductor]\label{lem:coins}
If positive integers $u,v$ are coprime, every integer
$t\geq(u-1)(v-1)$ equals $\alpha u+\beta v$ for some
$\alpha,\beta\in\Z_{\geq0}$.
\end{lemma}

\begin{proof}
Choose $\beta\in\{0,\ldots,u-1\}$ such that $\beta v\equiv t\pmod u$.
Then
$t-\beta v\geq (u-1)(v-1)-(u-1)v=-(u-1)>-u$.
Since it is a multiple of $u$, it is nonnegative.
\end{proof}

\begin{proof}[Proof of Theorem~\ref{thm:main}]
Put $c=g-1$. Start with $K_4$ and take $a$ Haj\'os joins with copies of
$K_4$, $b$ joins with copies of the five-rim wheel, and $d$ joins with
copies of $M(C_5)$, where $a,b,d$ are arbitrary nonnegative integers.
Lemma~\ref{lem:hajos} gives a 4-critical skeleton $H$ with
\[
|V(H)|=4+3a+5b+10d,\qquad |E(H)|=6+5a+9b+19d.
\]
Substitution using Lemma~\ref{lem:sub} therefore realizes every order
\begin{align*}
N&=4+6c+a(3+5c)+b(5+9c)+d(10+19c)\\
 &=6g-2+a(5g-2)+b(9g-4)+d(19g-9).
\end{align*}
Write $N_*=6g-2$, $A=5g-2$, $B=9g-4$, $C=19g-9$.
The key identity is
\begin{equation}\label{eq:gcd}
9A-5B=2.
\end{equation}

If $g$ is odd, $A,B$ are odd, so~\eqref{eq:gcd} implies
$\gcd(A,B)=1$. Lemma~\ref{lem:coins} with $d=0$ realizes every order
\[
N\geq N_*+(A-1)(B-1)=45g^2-46g+13.
\]

If $g$ is even, $A,B$ are even, $\gcd(A/2,B/2)=1$, and $C$ is odd.
Put $Q=(A/2-1)(B/2-1)$. For any remainder $t=N-N_*\geq C+2Q$,
choose $d=0$ when $t$ is even and $d=1$ when $t$ is odd.
Then $t-dC$ is even and at least $2Q$.
Lemma~\ref{lem:coins} represents $(t-dC)/2$ using $A/2,B/2$.
Consequently every order
\[
N\geq N_*+C+2Q=\frac{45}{2}g^2-8g+1
\]
is realized. Both sufficient thresholds are at most $45g^2$ for
$g\geq1$. Every realization retains the required vertex-criticality
and edge tolerance by Lemma~\ref{lem:sub}.
\end{proof}

\section{Scope}

The transfer theorem is a finite, elementary construction independent of
how the seed is obtained. Applying the asserted seed theorem of~\cite{KK}
closes the all-orders implication described in~\cite[Problem 1.1]{SS}.
The substitution lemma and Haj\'os operation are prior ingredients;
the arithmetic completion above supplies the additional argument.
The source search did not establish whether this completion has appeared
elsewhere. No explicit growth rate for $f_4(N)$ is claimed without a bound
on the order $g(r)$ of a seed. No Lean formalization or independent
re-proof of the cyclic and Chinese-remainder seed construction is included.

\begin{thebibliography}{9}
\bibitem{KK}
Liam Kruer and Jensen Kohlmeyer,
\emph{Vertex-critical four-chromatic graphs with no small critical edge set},
Theorem 1.1, source manuscript consulted 6 November 2026.
\url{https://github.com/JENW1N/complete-math-solutions/blob/main/papers/erdos-944.pdf}.

\bibitem{SS}
Skottov\'a and Steiner, arXiv:2508.08703,
especially Problem 1.1, Lemmas 2.1--2.2 and Theorem 2.3.
\url{https://arxiv.org/html/2508.08703}.

\bibitem{EP}
Erd\H{o}s Problems, Problem 944 and its history.
\url{https://www.erdosproblems.com/history/944}.
\end{thebibliography}
\end{document}
